\section{Discussion}\label{sec:discussion}

\subsection{What the results establish}
Four findings stand on consistent evidence. First, sharing light with the crop is the main lever on the electricity that tracker-based agrivoltaic arrays can produce under a yield floor. With static trackers the floor limits arrays to about 6\% ground coverage and an LER of 1.07. Any light-sharing control allows about three times the density, and the LER rises to 1.18--1.24 for hand-written rules. Second, the best simple rule follows crop phenology. It shares light across the middle of the season, around flowering and grain filling, and produces electricity at the beginning and the end. It needs only the crop's thermal time, which is observable on the farm, and it is insensitive to forecast errors, whereas the rule that reacts to heat and drought stress is worse in every climate and degrades with forecast error. Third, the qualitative result does not depend on the weakest part of the model. The simulated shade response is conservative against the field trials we could use (Section~\ref{sec:shade}); calibrating it to a meta-analysis or fitting it to the trials raises the allowed density, and the phenology rule stays the best method, with a gain over static arrays of 0.17, 0.19 and 0.30 LER under the conservative, calibrated and field-fitted responses. The absolute gains reported are therefore lower bounds with respect to those trials, not upper bounds, but the trials are few, cover four crops and one or two seasons each, and disagree with the meta-analysis for maize. Fourth, a controller learned through the differentiable simulator is better than the best rule we could write, by 0.03--0.05 in LER and about 10\% in electricity at equal retention, reproducibly over training seeds (range 0.002 in LER), three spatial splits, three canopy-cooling assumptions and four climates, and at 84--86\% of the held-out sites. This gain is small compared with the gap between static and dynamic control.

\subsection{What role does artificial intelligence play?}\label{sec:ai}
The results show where learning matters in this problem and where it does not, and we state this directly because the conclusions would otherwise be overstated. The largest effect, light sharing itself, comes from a physical design choice that a hand-written rule captures. The learned controller refines the timing of that rule and adds 0.03--0.05 LER; a reader who needs a simple, auditable controller can use the phenology rule and obtain about 85\% of the electricity gain of the learned controller over static arrays. The learned policy also has to be initialised from the rule, because training from random weights collapsed to full sharing (Section~\ref{sec:limits}), so we cannot claim that learning discovers the schedule on its own.

Three contributions are specific to the learning approach. First, differentiability is what made the evaluation possible: it allowed us to optimise a policy through whole 365-day seasons for hundreds of crop--site pairs and to screen design variants (density conditioning, floor conditioning, chance-constrained training, forecast noise) at a cost that black-box reinforcement learning \citep{gautron2022} would not allow. Second, the learned policy served as a test of hypotheses about the problem. A network with access to heat, drought and forecast inputs did not use them, and a network without them performed equally well, which supports the conclusion that the benefit comes from phenology and not from stress avoidance, a conclusion that rule design alone could not have reached. Third, one network trained at a single floor remains better than the phenology rule at floors between 0.80 and 0.95, which is a practical advantage for regulations with different floors (+0.033 to +0.037 LER, baseline climate). Conditioning the network on the floor and chance-constrained training gave no further improvement (Section~\ref{sec:extra}). These are, however, modest contributions. A stronger case for learning would need a setting in which the hand-written rule cannot be tuned to the answer, for instance control under uncertainty in the crop response to shade, where one policy has to remain safe across plausible responses, or per-site adaptation from on-farm measurements. We did not test these, and we regard them as the main route by which artificial intelligence could add value beyond what this paper shows.

\subsection{What the results do not establish}
The early hypothesis that stress-aware control, that is, shading the crop on heat and drought days, would be a main source of benefit under warming was not supported. The stress rule performs worse than the phenology rule in every climate, and a learned controller without stress inputs performs as well as one with them. A plausible reason is that in our crop model the yield gain from shading on a hot or dry day is smaller than the loss of light on that day, so stress information does not help. Whether this holds for crop models with more detailed canopy-temperature and heat-stress responses is an open question. Agrivoltaics does not protect staple production against warming in our simulations: light-sharing systems produce about 90\% of the open-field yield in every climate and follow the same decline in maize and rice (Fig.~\ref{fig:warming}).

\subsection{Food security in practice}
Holding the historical mean retention at 90\% does not hold it in each season or at each site. 70\% of SALS seasons meet the floor, the lower tail is slightly heavier than for the rules, and production-weighted retention is below 0.9 for soybean (0.888) and maize (0.898) and in several regions. Chance-constrained design raises season-level compliance to about 92\% at a cost of 0.065 in LER, and the ranking of the controllers is preserved. Whichever formulation is chosen, a per-site regulatory floor would need a region- and crop-specific design step. The upscaled potential (Section~\ref{sec:global}) shows the trade-off in magnitude. The 14\,PWh\,yr$^{-1}$ that 5\% of the staple cropland could supply need about 10\,TWp of modules, several times the present global fleet, and the best sites are rainfed maize and soybean areas where production-weighted retention is lowest.

\subsection{Limitations}\label{sec:limits}
\begin{itemize}
\item \emph{Simulation only.} All results come from a simulator. PV output was benchmarked against PVGIS (29 sites, mean bias $-4.9$\%), but the crop response to shade was only compared with a handful of field trials and one meta-analysis (Section~\ref{sec:shade}). The simulated response is close to proportional, is conservative with respect to the trials (rice, wheat, potato, maize) and does not reproduce the stronger losses of maize and legumes in \citet{laub2022}. The harvest-index sensitivities $h$ are assumptions, and the sensitivity analysis uses rule controllers on a subsample of 22 sites.
\item \emph{Canopy cooling.} The effect of shade on canopy temperature was represented by one coefficient $\kappa$, applied to a heat-stress function that was calibrated on air temperature \citep{siebert2014,webber2018}. The ranking of the controllers is robust to $\kappa$ between 0 and 2\,\textdegree C, but the absolute electricity changes by about $\pm8$\%. Nocturnal warming under modules and rain redistribution are not modelled.
\item \emph{Crop model.} SIMPLE has a fixed harvest index and linear radiation-use efficiency. More detailed models with grain-number formation and canopy energy balance may favour different schedules.
\item \emph{Climate scenarios.} Warming was applied as uniform temperature shifts with fixed precipitation, radiation and relative humidity, and with unchanged cultivars and sowing dates. It is not a bias-adjusted climate projection.
\item \emph{Initialisation and state information.} The network is initialised from the phenology rule, because training from random weights collapsed to full sharing. The gain of SALS is therefore the gain of fine-tuning that rule, and we cannot say whether a better optimiser or a richer policy class would find larger gains. The controller uses the simulated state of the crop and soil (thermal time, canopy interception, soil water) without measurement error, and the forecast errors we tested apply to weather only.
\item \emph{Economics and deployment.} There is no cost, market or grid analysis. The value of midday electricity, curtailment, storage, structural cost and farmer acceptance \citep{trommsdorff2021} will determine the realisable share. The average row pitch at the chosen densities (about 13\,m for 2.2-m modules) leaves room for machinery, but the densest designs (GCR up to 0.5) have a pitch of only 4.4\,m.
\end{itemize}
